\section{The Configuration Tax}
\label{sec:background}

\subsection{Configuration as glacial module state}
A production Earth-system or materials model is parameterized by a configuration
file---a Fortran \emph{namelist} of a few hundred variables that select solvers,
toggle physics, and fix grid and spectral dimensions. At startup the model executes
\lstinline[style=fortran]|READ(NML=...)| once, depositing each value into a field
of a derived type in a configuration module (in ICON, \texttt{mo\_run\_config},
\texttt{mo\_dynamics\_config}, \texttt{mo\_nonhydrostatic\_config}). Every consumer
then reaches the value through \lstinline[style=fortran]|USE module, ONLY: var|.
These variables are written exactly once, at initialization, and only read during
the thousands of timesteps that follow: they are \emph{glacial}~\cite{glacial}. They
are also exactly the values a compiler cannot see---read behind an I/O barrier,
held in module-level state, and consumed in a compute phase the compiler does not
distinguish from initialization.

\subsection{Why \texttt{-O3} and LTO do not help}
\label{sec:whycompilers}
It is tempting to assume that whole-program optimization recovers these constants.
To test this precisely we take the NAS~LU benchmark as a single translation unit
whose \texttt{COMMON} blocks are rewritten to module variables---the same shape as a
configuration module---and ask whether the LLVM pipeline can fold a representative
grid bound (\texttt{nz}) into the dozens of loops that use it.
\Cref{tab:lu} reports the result across three settings (flang-new-21,
LLVM~21).\footnote{Full reproduction (HLFIR + LLVM IR at \texttt{-O0/-O2/-O3} and the
whole-program LTO test): \texttt{experiments/lu\_globals/reproduce.sh}.}

\begin{table}[t]
  \centering
  \caption{\textbf{[PROVISIONAL --- experiment under revision, see PLAN.md.]} Folding
  a configuration-style module global (\texttt{nz}, a loop bound) in NAS~LU. SCCP never
  participates (an SSA-value pass; globals are memory). Rows A/B are solid; row C is
  currently only a presence-gated partial fold (the injection was a runtime store, not
  a source constant) and must be redone.}
  \label{tab:lu}
  \small
  \begin{tabular}{@{}lll@{}}
    \toprule
    Setting & Linkage & \texttt{nz} bound folded? \\
    \midrule
    A. as-is, \texttt{-O3} (value never set in TU)        & external & \textbf{no} --- value absent \\
    B. runtime store, single-TU \texttt{-O3}              & external & \textbf{no} --- external linkage \\
    C. runtime store, whole-program LTO + internalize     & demoted to \texttt{i1} & \textbf{partial} --- presence-gated \\
    \bottomrule
  \end{tabular}
\end{table}

\td{TODO/CORRECTION (skeptical review): row C does NOT cleanly fold \texttt{nz}.
GlobalOpt demoted the global to a one-bit "has-init-run" flag and rematerialized the
bound as \texttt{select i1 \%flag, i32 33, i32 0} (6 sites; 21 others got literal 33),
because the value was injected as a RUNTIME STORE, not a source constant. REDO injecting
a Fortran \texttt{parameter} / constant initializer (what config-prop produces) for a
clean fold; add a row with a real \texttt{READ(NML=...)} so row A is literally the I/O
barrier; tie the metric to "bound is a compile-time constant", not load-counts.}

Rows A and B isolate two real barriers: \emph{value presence} (A --- with no constant
in the program nothing can be folded) and \emph{linkage} (B --- an external-linkage
module/\texttt{COMMON} global is not foldable single-TU; a unit must assume another may
write it). \td{QUESTION: is the third regime (a true source-constant injection in a
closed program) a CLEAN fold? Until the redo, claim only that config-prop supplies a
source-level constant the pipeline otherwise lacks; do not assert "LTO folds it."}
Production configuration state fails the value-presence and linkage conditions; whether
the residual is "structural" in the strong sense needs the redone row C.

\subsection{What deployment-time specialization adds}
\Cref{tab:lu} also delimits what is \emph{not} novel: once a scalar value is supplied
\emph{as a source constant} to a closed program, the standard cascade---constant
propagation, dead-code elimination, fusion---does the rest (subject to the row-C redo). Our contribution for
that case is the bridge of \Cref{sec:recovery}: recovering the value statically from
the configuration file, the one input the compiler structurally cannot obtain. The
contributions that go \emph{beyond} a supplied scalar---narrowing index storage
from a recovered interval, and devirtualizing from a recovered type-set
(\Cref{sec:spec})---are optimizations the right column of \Cref{tab:lu} does not
perform at any setting, and are where the bulk of our speedup originates.

\td{DEEP ARG A (PLAN.md 2.A) --- phase-ordering / config-prop reparameterizes the
INPUT. Toolchain $B=(T_n\circ\dots\circ T_1)(S)$ ends in LTO; the config value
$v=F(\mathrm{var})$ is not derivable from $S$ at any phase, so the terminal LTO sink
cannot fold it. Config-prop is not a phase: it computes
$\mathrm{compile}(\mathrm{inject}(S,F))$. LU row C is this with inject done by hand;
row A is $\mathrm{compile}(S)$. "Why not LTO" fails on phase-ordering, not capability.}
\td{DEEP ARG D (PLAN.md 2.D) --- the "isn't this just recompile-per-config?"
three-layer floor: (1) UNAVAILABLE (values are runtime reads into mutable external
module state; baking in needs the glacial may-write proof = the contribution);
(2) INFEASIBLE (full per-config rebuild of a 1M-line model at every site; covering set
amortizes); (3) INSUFFICIENT (a scalar narrows no index type, devirtualizes no
dispatch). Consider a short Discussion subsection.}
